% ============================================================
% LuaLaTeX preamble for Japanese article
% Japanese: Gothic
% English: Sans-serif
% ============================================================

% --------------------
% Page layout
% --------------------
\usepackage[
  top=10mm,
  bottom=20mm,
  left=15mm,
  right=15mm
]{geometry}

% --------------------
% Japanese / English fonts
% --------------------
\usepackage{luatexja}
\usepackage[haranoaji,deluxe]{luatexja-preset}

\usepackage{fontspec}
\usepackage{unicode-math}

% English: sans-serif
\setmainfont{TeX Gyre Heros}
\setsansfont{TeX Gyre Heros}
\setmonofont{Latin Modern Mono}

% Math
\setmathfont{Libertinus Math}

% Japanese: Gothic
\renewcommand{\kanjifamilydefault}{\gtdefault}

% % Use sans-serif as the default Latin family
% \renewcommand{\familydefault}{\sfdefault}

% --------------------
% Basic packages
% --------------------
\usepackage{xcolor}
\usepackage{graphicx}
\usepackage{float}

% --------------------
% Tables
% --------------------
\usepackage{booktabs}
\usepackage{array}
\usepackage{tabularx}
\usepackage{longtable}
\usepackage{multirow}
\usepackage{makecell}

% --------------------
% Lists
% --------------------
\usepackage{enumitem}
\setlist[itemize]{leftmargin=2em}
\setlist[enumerate]{leftmargin=2em}

% --------------------
% Header / footer
% --------------------
\usepackage{fancyhdr}
\pagestyle{fancy}
\fancyhf{}
\fancyfoot[C]{\textcolor{white!50!black}{\textsf{LLM-Driven Utility Ellicitation}}}
\renewcommand{\headrulewidth}{0pt}
\renewcommand{\footrulewidth}{0pt}

% --------------------
% Icons
% --------------------
\usepackage{fontawesome5}

% --------------------
% Hyperlinks
% --------------------
\usepackage{hyperref}

\hypersetup{
  unicode=true,
  colorlinks=true,
  linkcolor=teal!80!black,
  citecolor=teal!80!black,
  urlcolor=teal!80!black
}

% --------------------
% Math packages
% --------------------
\usepackage{amsmath}
\usepackage{mathtools}
\usepackage{bm}

% --------------------
% Theorem environments
% --------------------
\usepackage{amsthm}

\theoremstyle{plain}
\newtheorem{thm}{Theorem}[section]
\newtheorem{proposition}[thm]{Proposition}
\newtheorem{lemma}[thm]{Lemma}
\newtheorem{corollary}[thm]{Corollary}

\theoremstyle{definition}
\newtheorem{definition}[thm]{Definition}
\newtheorem{assumption}[thm]{Assumption}
\newtheorem{exam}[thm]{Example}

\theoremstyle{remark}
\newtheorem{remark}[thm]{Remark}

% --------------------
% Operators and notation
% --------------------
\DeclareMathOperator{\Avar}{Avar}
\DeclareMathOperator{\Var}{Var}
\DeclareMathOperator{\varop}{var}
\DeclareMathOperator{\Cov}{Cov}
\DeclareMathOperator{\cov}{cov}
\DeclareMathOperator{\rank}{rank}
\DeclareMathOperator{\diag}{diag}
\DeclareMathOperator{\MSE}{MSE}
\DeclareMathOperator{\IMSE}{IMSE}
\DeclareMathOperator{\bias}{bias}

\DeclareMathOperator*{\argmax}{arg\,max}
\DeclareMathOperator*{\argmin}{arg\,min}

\newcommand{\E}{\mathbb{E}}
\newcommand{\R}{\mathbb{R}}
\newcommand{\Hbb}{\mathbb{H}}
\newcommand{\Pbb}{\mathbb{P}}
\newcommand{\Sbb}{\mathbb{S}}

\newcommand{\parrow}{\xrightarrow{p}}
\newcommand{\darrow}{\xrightarrow{d}}
\newcommand{\asarrow}{\xrightarrow{a.s.}}

\newcommand{\plim}{\operatorname*{plim}}
\newcommand{\iid}{\text{i.i.d.}}
\newcommand{\Normal}{\text{Normal}}
\newcommand{\Unif}{\text{Uniform}}
\newcommand{\CI}{\text{CI}}

\newcommand{\what}[1]{\widehat{#1}}

% --------------------
% Paragraph spacing
% --------------------
\setlength{\parindent}{1em}
\setlength{\parskip}{0.3em}